
\chapter{Платёжный шлюз}\label{ch:gateway}

Один запрос продавца~--- \texttt{POST} к~\texttt{/v1/payment\_intents}
на~оплату заказа. Снаружи это обычный JSON (JavaScript Object Notation)
поверх HTTPS. Внутри запрос проходит через последовательность решений,
и~ошибка на~любом шаге меняет риск, конверсию или PCI~scope~--- область
оценки соответствия PCI~DSS (Payment Card Industry Data Security
Standard)~\cite{pci-dss-4}.

\section{Устройство шлюза}\label{sec:gw-anatomy}

\highlight{Платёжный шлюз (payment gateway) превращает один запрос продавца
в~корректную, наблюдаемую авторизацию}. Разграничение
шлюза, процессора и~оркестратора как ролей экосистемы дано
в~разделе~\ref{sec:ecosystem-participants}.

Внутри шлюз делится на~модуль приёма запросов программного
интерфейса (API, application programming interface), модуль оценки риска,
модуль маршрутизации, изолированное хранилище карточных данных и~токенов
(vault) и~модуль авторизации, отправляющий запрос банку-эквайеру
или напрямую карточной сети. Порядок их вызова определяет
и~архитектурные свойства системы, и~её PCI scope.

Продавец отправляет HTTPS-запрос на~API шлюза, например,
\texttt{POST} к~\texttt{/v1/payment\_intents} либо
\texttt{/v1/charges} в~более старой версии API, с~JSON-телом,
содержащим сумму, валюту, данные карты или токен и~метаданные
заказа. Модуль приёма валидирует сообщение, проверяет
аутентификацию по~API-ключу или токену протокола делегированной
авторизации OAuth~2.0~\cite{rfc-6749} и~идемпотентность,
подробнее в~разделе~\ref{sec:gw-reliability}, затем передаёт
нормализованное внутреннее представление транзакции модулю оценки риска.

Модуль оценки риска применяет к~транзакции набор правил и~моделей:
проверки частоты повторов (velocity checks), сопоставление
геолокации, отпечаток устройства (device fingerprint, см.
главу~\ref{ch:fraud}), а~в~зрелых системах ещё и~скоринговую модель
вероятности мошенничества. Решений три: пропустить транзакцию
дальше, отправить её на~дополнительную проверку по~протоколу
подтверждения держателя карты 3-D~Secure или отклонить. Устройство
систем антифрода разобрано в~главе~\ref{ch:fraud}. Проверка риска
работает \emph{до} авторизации: подозрительную транзакцию дешевле
и~проще остановить заранее, чем потом разбирать чарджбэк.

Если транзакция прошла оценку риска, она попадает в~модуль
маршрутизации, который определяет, через какого эквайера
или процессор отправить авторизационный запрос
(см.~раздел~\ref{sec:gw-routing}).

Если продавец передал полный номер карты (PAN, Primary Account Number),
шлюз заменяет его токеном из~хранилища. PAN сохраняется
в~зашифрованном хранилище (vault), а~дальше по~внутреннему конвейеру
и~в~ответе продавцу циркулирует только токен. Это сужает CDE
(cardholder data environment) до~одного компонента и~радикально
сокращает PCI~scope всей системы~--- механизм подробно разобран
в~разделе~\ref{sec:pci-scope-reduction}. Если продавец с~самого начала
прислал сетевой токен (DPAN, device PAN~--- токенизированный
аналог реального PAN), хранилище может не~участвовать, но~шлюз
всё равно проверяет ограничения домена токена и~при необходимости
запрашивает у~провайдера сетевой токенизации (TSP, Token Service
Provider) актуальную одноразовую криптограмму транзакции
для~конкретной операции~\cite{emvco-payment-tokenisation-guide}
(см.~также раздел~\ref{sec:token-anatomy}).

Модуль авторизации формирует сообщение по~международному стандарту
карточных операций ISO~8583 (см.~главу~\ref{ch:iso8583}) или его
эквивалент в~протоколах API новых эквайеров, отправляет запрос
по~защищённому каналу банку-эквайеру, получает ответ от~эмитента
через карточную сеть, преобразует код ответа в~формат API шлюза
и~возвращает результат продавцу.

Открытая реализация всех описанных компонентов в~одной кодовой
базе~--- Hyperswitch от~Juspay~\cite{hyperswitch-juspay} (Rust,
Apache~2.0): унифицированные коннекторы к~120+~процессорам, модуль
маршрутизации с~правилами повтора, vault для~токенов и~карточных
данных, соответствующий PCI~DSS. Та же
кодовая база работает у~самого Juspay в~промышленной эксплуатации.

\begin{figure}[htbp]
  \centering
  \includefiguresvg[width=\linewidth]{assets/figures/ch33-payment-gateway/gateway-components}
  \caption{Компоненты платёжного шлюза.}\label{fig:gateway-components}
\end{figure}

\section{Оркестрация поверх шлюза}\label{sec:gw-orchestration-layer}

Шлюз, процессор и~оркестратор~--- три разных архитектурных уровня;
различие становится продуктовым, когда ТСП подключает
второго эквайера. Разбор ролей экосистемы дан
в~разделе~\ref{sec:ecosystem-participants}.

\textbf{Шлюз} (см.~раздел~\ref{sec:gw-anatomy})~--- протокольная интеграция
с~одним эквайером или картной сетью. Преобразует JSON-запрос в~ISO~8583
(или REST-эквивалент современных процессоров), держит соединение,
токенизирует PAN, валидирует ответ.

\textbf{Процессор} обеспечивает сетевой обмен между шлюзом и~эмитентом
через карточную сеть: маршрутизирует сообщение в~нужную сеть
(Visa, Mastercard, Мир), при необходимости подменяет недоступного
эмитента собственным ответом по~правилам \textit{stand-in
processing}~\cite{mc-switching-services,nspk-mir-rules}. У~одного провайдера
процессор и~шлюз нередко совмещены в~один узел.

\textbf{Оркестратор}~--- уровень над несколькими PSP и~шлюзами.
Продукт обращается к~единому API, оркестратор разводит вызовы
по~десяткам шлюзов: маршрутизирует транзакцию между провайдерами
по~правилам стоимости, доли одобрений и~регионального профиля;
реализует каскадный повтор при отказе одного из~маршрутов
(см.~раздел~\ref{sec:gw-routing}); держит единый vault, унифицирует форму
оплаты (\textit{checkout}) поверх разных интеграций. Маршрутизация
и~оценка риска, описанные внутри одного шлюза, на~уровне оркестратора
повторяются между шлюзами.

Оркестратор окупается при типовых нагрузках:

\begin{itemize}
  \item ТСП работает с~двумя и~более PSP одновременно~---
    ради устойчивости, ради арбитража стоимости или ради
    разнесения трафика по~сетям;
  \item региональная экспансия, при которой каждому рынку
    нужен локальный эквайер с~собственным договором и~процессингом;
  \item маркетплейс, где продавцы хотят выбирать собственного PSP
    или подключать унаследованный договор эквайринга;
  \item стратегический уход от~привязки к~одному вендору, когда
    переключение PSP должно занимать дни, а~не~кварталы.
\end{itemize}

Primer~\cite{primer-orchestration} предлагает конфигурацию
маршрутизации без программирования (\textit{no-code}) поверх большого
набора подключений; типичные вертикали~--- travel, retail, gaming,
marketplaces. Spreedly~\cite{spreedly-orchestration}~--- открытая
независимая от~PSP (\textit{PSP-agnostic}) платформа с~120+~интеграциями
и~универсальным vault уровня PCI~Level~1.
Gr4vy~\cite{gr4vy-orchestration}~--- однопользовательская
(\textit{single-tenant}) платформа с~облачной (\textit{cloud-native}) архитектурой,
устраняющая привязку к~одному вендору
(\textit{vendor lock-in}). Hyperswitch~\cite{hyperswitch-juspay}
от~Juspay~--- открытый исходный код на~Rust под Apache~2.0,
120+~подключаемых процессоров и~vault, соответствующий PCI~DSS;
сам Juspay использует ту же
кодовую базу в~промышленной эксплуатации и~в~FY25 обработал
свыше 300\,млн~транзакций в~сутки при~годовом TPV около~\$1\,трлн.

В~2025~году рынок payment orchestration оценивается
в~\$3{,}5\,млрд с~прогнозом \$18\,млрд к~2034~году
и~CAGR около~20\,\% (отчёт Intel Market
Research)~\cite{intel-market-payment-orchestration}.

\highlight{Оркестратор появляется в~тот момент, когда у~ТСП
одновременно работают несколько шлюзов с~разной экономикой;
на~архитектурной диаграмме он занимает отдельное место между
приложением продавца и~всем парком PSP}.

\section{Маршрутизация запросов}\label{sec:gw-routing}

При нескольких эквайерах шлюз решает, куда
отправить авторизационный запрос, балансируя стоимость,
долю одобрений и~время ответа.

\subsection*{Критерии оптимизации}

\textbf{Стоимость.} Тарифы у~эквайеров разные, а~interchange fee
зависит от~типа карты, четырёхзначного кода категории ТСП
(MCC, Merchant Category Code), региона эмитента и~десятков других
параметров. Модуль маршрутизации выбирает эквайера с~минимальной
стоимостью для~конкретной комбинации условий~--- например,
направляет внутрироссийскую транзакцию по~карте Мир через
эквайера с~более выгодным тарифом Национальной системы платёжных
карт (АО НСПК), а~международную~--- через другого, у~которого ниже
валютная наценка.

\textbf{Доля одобрений.} На~одних и~тех же диапазонах BIN (Bank
Identification Number, идентификатор банка-эмитента в~составе PAN)
эквайеры показывают разную долю одобрений. Причины бывают техническими
или коммерческими: качество подключения к~сети, скорость ответа,
репутация ТСП у~конкретного эквайера, локальный или международный
маршрут. Накопив статистику по~BIN, стране и~сумме, модуль
маршрутизации направляет транзакцию туда, где вероятность одобрения
выше. Для~крупных объёмов даже небольшая разница в~качестве маршрута
даёт заметное изменение выручки.

\textbf{Латентность.} Продавцу важно, чтобы ответ пришёл быстро.
Если один эквайер стабильно отвечает заметно быстрее другого,
модуль маршрутизации при прочих равных предпочтёт первый: каждая
лишняя секунда в~сценарии оплаты снижает конверсию.

Эти три критерия противоречат друг другу. У~самого дешёвого эквайера
доля одобрений может быть хуже, а~самый быстрый~--- оказаться
самым дорогим. Модуль маршрутизации ищет оптимум в~ограничениях,
заданных продавцом: \enquote{минимизировать стоимость при доле
  одобрений не~ниже X\% и~p95~--- 95-м процентиле распределения~---
латентности не~более Y~мс}. Адаптивное распределение трафика между
маршрутами моделируют как задачу о~\enquote{многоруком бандите}
(multi-armed bandit).

\needspace{5\baselineskip}
\subsection*{Каскадная маршрутизация}

\enquote{Мягкий отказ} (soft decline)~--- это отказ, формально устранимый:
повторная попытка по~другому маршруту или после дополнительной
аутентификации может завершиться одобрением. Если первый маршрут
вернул такой отказ, шлюз отправляет транзакцию второму эквайеру. Сам термин
у~схем имеет более узкий нормативный смысл. Mastercard, например,
требует мягкий отказ в~странах со~строгой аутентификацией клиента
(SCA, Strong Customer Authentication) только тогда, когда
операции типа CNP (Card Not Present, операции без физического
предъявления карты) требуется аутентификация, но~она
отсутствует,~--- а~не~как общий синоним любого отказа~\cite{mc-tpr}.
Общая логика каскада от~этого не~меняется: отказ одного маршрута
ещё не~означает, что транзакция в~принципе невозможна. Другой эквайер
может получить иной ответ от~того же эмитента, особенно когда
его авторизационный запрос отличается: другой BIN эквайера,
другой MCC, другой идентификатор терминала (Terminal~ID, TID;
в~паре с~идентификатором ТСП, MID).

\importantnote{%
  Каскад и~обычный повтор решают разные задачи: каскад идёт к~\emph{другому}
  эквайеру, повтор~--- \emph{тому же} спустя время. Путаница порождает
  дубли: эмитент видит два запроса и~одобряет оба.
}

Каскад сам по~себе не~создаёт дубля на~стороне первого эквайера,
но~даёт другой риск: задержавшийся ответ от~первого маршрута
может прийти после того, как второй успел получить одобрение.
Шлюз немедленно отправляет отмену первой авторизации, если та
вернулась с~опозданием.

Каскад работает только для~части отказов, которые эквайер
или схема считают потенциально устранимыми. К~явно неверным
реквизитам~--- \texttt{14}~Invalid Card Number или
\texttt{54}~Expired Card~--- это не~относится~\cite{mc-tpr}.
Сами схемы ограничивают избыточные повторные попытки; конкретные
лимиты зависят от~сценария и~правил конкретной схемы~\cite{mc-tpr}.
Модуль маршрутизации обязан различать эти два класса отказов
по~коду ответа и~включать каскад только для~устранимых;
рабочая классификация кодов на~каскадные, ретраемые и~финальные
разобрана в~разделе~\ref{sec:decline-response-codes}.

Разница каскада, повтора и~финального отказа становится явной,
если выразить правило выбора чистой функцией. Вход~--- история
попыток и~список доступных эквайеров; выход~--- одно из~трёх
решений. Таблицы кодов в~примере ниже иллюстративны: в~промышленной
системе их задают правила схемы и~спецификация процессинга.

\codefile{Python}{samples/ch33-cascade-retry/retry.py}

\section{Интеграционные модели}\label{sec:gw-integration}

Способ подключения продавца к~шлюзу определяет, как
делится ответственность за~карточные данные и~каков
PCI scope продавца.

\subsection*{Размещённая платёжная страница (перенаправление)}

В~модели размещённой платёжной страницы (hosted payment page, HPP)
продавец инициирует платёжную сессию через API шлюза
и~получает URL, на~который перенаправляет покупателя. Покупатель
вводит данные карты на~странице шлюза, шлюз обрабатывает транзакцию
и~возвращает покупателя на~сайт продавца с~результатом. Параллельно
шлюз отправляет продавцу асинхронное уведомление с~финальным
статусом.

HPP сводит PCI~scope к~минимуму. Карточные данные не~проходят
через серверы продавца, поэтому он обычно попадает в~модель SAQ~A
(Self-Assessment Questionnaire~A, упрощённая анкета самооценки
PCI~DSS); более широкие сценарии, где элементы платёжной страницы
частично исходят от~самого продавца, к~нему не~применяются~\cite{pci-faq-1291,pci-faq-1293}.
Для~небольших ТСП без выделенной команды безопасности это часто
единственный реалистичный вариант.

Недостаток~--- потеря контроля над пользовательским сценарием.
Перенаправление прерывает поток оплаты, страница шлюза может
выбиваться из~визуального языка сайта, обратный переход
после оплаты не~всегда проходит гладко. При такой схеме сайт
продавца остаётся в~PCI~scope: PCI~SSC (PCI Security
Standards Council) отдельно уточняет, что redirect mechanism
оставляет сайт продавца в~области оценки~\cite{pci-faq-1332}.

\subsection*{Встроенная форма (iFrame / JavaScript SDK)}

Продавец встраивает в~свою страницу платёжную форму шлюза~---
встроенный HTML-фрейм (iFrame) или готовый JavaScript-компонент
из~набора инструментов разработчика (SDK, software development kit).
Визуально форма выглядит как часть сайта продавца, но~технически
данные карты вводятся внутри iFrame, контролируемого шлюзом,
и~отправляются напрямую на~серверы шлюза, минуя серверы продавца.

Модель сохраняет больше контроля над интерфейсом~--- продавец
стилизует форму, сохраняет цельный сценарий оплаты и~попадает
в~SAQ~A или SAQ~A-EP. Разница зависит
от~деталей реализации и~от~того, проходят ли карточные данные
через JavaScript продавца хотя бы транзитом~\cite{pci-faq-1291,pci-faq-1293}.

\importantnote{%
  PCI~DSS v4.0.1 ставит страницы оплаты под отдельный
  контроль~--- требования~6.4.3 и~11.6.1. В~январе~2025~года PCI~SSC
  обновил SAQ~A, заменив их упрощённым критерием: продавец
  подтверждает, что сайт не~подвержен атакам через вредоносные скрипты (FAQ~1588%
  ~\cite{pci-dss-4,pci-faq-1588}). Даже при вводе данных в~iFrame
  шлюза вредоносный скрипт продавца может перехватить ввод оверлеем
  или keylogger до~того, как данные дойдут до~iFrame, поэтому
  при встроенной форме продавцу нужно доказуемо контролировать
  скрипты на~payment page.
}

\subsection*{Прямая интеграция через API (server-to-server)}

Продавец собирает данные карты на~своей стороне через
собственную форму и~отправляет их на~API шлюза
в~серверном запросе. Это даёт полный контроль над
интерфейсом и~произвольную предварительную логику:
проверку адреса, предавторизацию, разделение платежа
между несколькими получателями. Карточные данные проходят
через серверы продавца, и~сценарий выпадает из~льготных моделей
SAQ~A / A-EP~\cite{pci-faq-1291,pci-faq-1293,pci-saq}, существенно
расширяя PCI scope.

Этот путь выбирают крупные ТСП и~платёжные платформы,
сертифицированные по~PCI~DSS Level~1, которым прямой
контроль над данными даёт собственное защищённое хранилище токенов
и~оптимизацию логики повторов на~уровне PAN.

\practicenote{%
  Начинать разумно с~iFrame или SDK~--- приемлемый пользовательский
  опыт при минимальном PCI scope. На~server-to-server переходят,
  когда ограничения встроенной формы блокируют
  конкретный сценарий, например, \textit{card-on-file} (CoF~---
  сохранённые в~продукте реквизиты карты для~повторных списаний)
  с~прозрачной для~покупателя токенизацией или каскадную
  маршрутизацию на~уровне PAN.
}

\section{Надёжность и~наблюдаемость}\label{sec:gw-reliability}

Надёжность шлюза определяется тем, как он справляется
с~неопределённостью: ответ запаздывает, внешний вызов
теряется, продавец нажимает \enquote{повторить}.

\subsection*{Идемпотентность}

Повторный вызов с~тем же набором параметров обязан давать тот же
результат, что и~первый, без побочных эффектов.
Если продавец отправил запрос на~авторизацию, не~получил ответа
из-за таймаута или сетевого сбоя и~повторил его с~тем же ключом
идемпотентности (idempotency key~--- клиентский идентификатор
операции, по~которому шлюз дедуплицирует повторные запросы),
шлюз возвращает результат первой попытки; вторая транзакция
не~создаётся~\cite{stripe-idempotency,adyen-api-idempotency}.

Для~этого шлюз хранит сопоставление между
\path{idempotency_key} и~\path{transaction_id}
в~хранилище со~строгой консистентностью~--- обычно
в~реляционной базе с~уникальным индексом по~ключу.
Срок жизни ключа должен покрывать весь период, когда
ещё возможны безопасные повторы одной и~той же операции,~---
от~сетевого сбоя сразу после ответа до~позднего повторного
запроса клиента.\footnote{Защита от~двойной обработки, встроенная
в~саму архитектуру слоя записи и~не~реализуемая в~каждом обработчике
отдельно, описана как сквозной принцип в~инженерном обзоре
собственного реестра Stripe~\cite{stripe-ledger-blog}: каждое
исправление оформляется новой проводкой со~ссылкой на~исходную,
а~разовые проверки \enquote{уже ли существует} не~используются.}

\importantnote{%
  Идемпотентность на~уровне шлюза не~заменяет идемпотентность
  на~уровне эквайера. Если шлюз отправил авторизацию,
  не~получил ответа и~отправил повторно, эквайер может
  обработать оба запроса. Шлюз либо использует поля данных
  ISO~8583 (data elements, DE) DE~11 (STAN, System Trace Audit
  Number~--- внутренний номер транзакции в~коротком интервале)
  и~DE~37 (RRN, Retrieval Reference Number~--- идентификатор
  для~последующих обращений к~транзакции) для~дедупликации
  эквайером, либо реализует сценарий \enquote{authorize-then-check}
  (\enquote{авторизуй-затем-проверь})~--- при таймауте отправляет
  не~повтор, а~запрос статуса (status~inquiry) у~эквайера, когда
  такой механизм предусмотрен протоколом, и~только при статусе
  \enquote{не~найдено} повторяет авторизацию.
}

\subsection*{Обработка таймаута}

Таймаут опасен тем, что отсутствие ответа~--- это \emph{не}
отказ. Если шлюз отправил авторизационный запрос эквайеру
и~не~получил ответа в~течение порога, возможны как минимум
три исхода: (а)~запрос не~дошёл; (б)~запрос дошёл,
эмитент одобрил транзакцию, но~ответ потерялся на~обратном пути;
(в)~запрос дошёл, эмитент отказал, ответ тоже потерялся.
В~момент таймаута шлюз не~различает эти случаи.

Наихудший исход~--- (б): деньги заблокированы на~карте
держателя, эмитент считает транзакцию одобренной, продавец
не~знает об~этом и~может повторить платёж, создав двойное списание.
Карточные правила решают эту неопределённость отменой.
Mastercard требует \textit{automatic reversal}~--- автоматическую
отмену авторизации при неполучении ответа на~\textit{authorization
request}~\cite{mc-tpr}. После таймаута шлюз сначала разрешает
неопределённое состояние и~только затем решает, что возвращать
продавцу.

\subsection*{Сбой шлюза между приёмом ответа и~записью}

Ответ пришёл, эмитент одобрил транзакцию, но~шлюз упал
до~того, как сохранил результат. Состояние транзакции
рассинхронизируется между шлюзом и~эмитентом:
эмитент удерживает деньги на~карте
держателя, а~шлюз после перезапуска находит запись в~статусе
\texttt{sent\_to\_acquirer} без сведений об~исходе.

Если шлюз просто повторит авторизацию, эмитент получит два
удержания на~одну сумму~--- двойное списание. Если шлюз вернёт
продавцу ошибку, деньги останутся заблокированными без
товара. Шлюз ведёт двухстатусный жизненный цикл
записи (\texttt{pending}~$\to$~\texttt{final}), который иногда
по~аналогии называют \enquote{двухфазной фиксацией}; с~распределённым
двухфазным коммитом (2PC) из~теории СУБД его роднит только идея
двух стадий~--- координации между несколькими узлами здесь нет.
До~отправки запроса эквайеру шлюз создаёт запись
\texttt{pending} с~idempotency key и~STAN. После перезапуска
фоновый процесс выбирает все \texttt{pending}-записи, по~которым
прошло время разумной задержки ответа эквайера, и~выполняет
запрос статуса (status~inquiry) к~эквайеру по~STAN и~RRN. Порог
здесь задаёт не~сетевой таймаут авторизации (он измеряется
десятками секунд, см.~ниже), а~время, после которого ответ
по~исходному каналу уже не~ожидается. Если эквайер подтверждает
одобрение, шлюз переводит запись в~\texttt{approved} и~ставит вебхук
(\textit{webhook}) в~очередь. Если эквайер не~знает о~транзакции, шлюз записывает
\texttt{failed} и~освобождает продавца для~повторной попытки.
Когда запрос статуса недоступен (не~все эквайеры его
реализуют), шлюз отправляет отмену, гарантированно
снимая холд, и~возвращает продавцу ошибку с~рекомендацией
повторить~\cite{mc-tpr}.

Без этого механизма после каждого перезапуска шлюза остаются
зависшие холды: держатели карт видят их как
списания, продавцы~--- как потерянные заказы. На~масштабе
тысяч транзакций в~минуту даже секундный сбой создаёт десятки
расхождений, разбор которых вручную через эквайера занимает дни.

\subsection*{Дедупликация}

Дедупликация~--- защита от~двойной обработки одной и~той же
транзакции~--- работает на~нескольких уровнях. На~уровне API она опирается
на~idempotency key. На~уровне внутреннего конвейера~---
на~уникальность комбинации STAN + RRN + amount + timestamp,
которая отлавливает дубли, проскочившие мимо идемпотентности:
например, при каскадной маршрутизации, когда два эквайера
получают запрос с~разными STAN, но~за~одну и~ту же покупку.
На~уровне денежных состояний дедупликацию обеспечивает
внутренний реестр обязательств из~главы~\ref{ch:internal-ledger},
а~на~уровне внешних файлов и~банковской выписки~---
сверка из~главы~\ref{ch:reconciliation}.

\subsection*{Вебхуки и~гарантия доставки}

Платёжный поток редко заканчивается в~момент синхронного
ответа API. Результат транзакции возвращается продавцу
двумя путями: синхронно, в~ответе на~запрос, и~асинхронно,
через обратный вызов по~HTTP (вебхук) на~указанный URL.
Синхронный ответ может не~дойти; кроме того, после первого
ответа появляются новые внешние события: завершился 3DS
challenge, пришло подтверждение, сработал возврат, поступило
позднее обновление статуса от~процессора.

Публичные API современных платёжных сервис-провайдеров (PSP,
payment service provider)~--- от~международных Stripe и~Adyen
до~местных ЮKassa, CloudPayments, T-Бизнес и~Сбер для~бизнеса~---
строят доставку вебхуков по~принципу
\enquote{как минимум один раз}. Эндпоинт должен быть доступен
по~HTTPS, система ждёт ответ 2xx, при неуспехе повторяет
доставку, а~продавец проверяет подпись события, прежде чем
доверять содержимому запроса. У~Stripe подпись
передаётся в~HTTP-заголовке \texttt{Stripe-Signature: t=...,v1=...}, % chktex 11
где \texttt{t}~--- timestamp подписи, а~\texttt{v1}~---
HMAC-SHA256 (Hash-based Message Authentication Code на~базе
SHA-256~\cite{rfc-2104}) от~конкатенации timestamp, точки
и~тела запроса с~использованием секрета подписи
вебхуков~\cite{stripe-webhooks}.\footnote{Правила подписи и~формат
вебхуков у~разных PSP различаются. У~одних подпись HMAC-SHA256
в~HTTP-заголовке, у~других~--- токен JWT (JSON Web
Token~\cite{rfc-7519}) или собственная схема;
список повторов и~таймаутов тоже разный. Перед интеграцией
с~конкретным PSP сверяйтесь с~его документацией; перенос ожиданий
от~Stripe на~любую другую интеграцию ведёт к~ошибкам.}

\practicenote{%
  Обработчик вебхуков на~стороне продавца обязан быть идемпотентным:
  до~изменения состояния заказа нужно проверить, не~обработан ли уже
  этот \path{event_id}.
}

\subsection*{Наблюдаемость}

При тысячах транзакций в~минуту метрики обязаны фиксировать деградацию
раньше, чем она затронет значительную долю трафика.

Доля одобрений~--- процент одобренных транзакций
от~общего числа попыток, с~разбивкой по~эквайеру, диапазону BIN,
стране эмитента, MCC и~сумме. Падение доли одобрений
за~короткий интервал~--- повод к~немедленному расследованию:
причиной бывает как сбой у~эквайера, так и~массовая атака
скомпрометированными картами.

Время ответа отслеживается через p50, p95 и~p99~--- 50-й, 95-й
и~99-й процентили распределения времени обработки запроса,
измеренного от~получения HTTPS-запроса продавца до~отправки
ответа. Неожиданный рост p99 указывает
на~сбои аппаратного криптомодуля (HSM, hardware security module,
см.~главу~\ref{ch:crypto}), медленный ответ эквайера или перегрузку
собственного модуля маршрутизации.

Доля системных ошибок включает транзакции,
завершившиеся таймаутом, сетевой или внутренней ошибкой сервера,
а~не~бизнес-ответом вроде \texttt{approved} или \texttt{declined}.
Рост этой доли на~фоне обычного трафика~--- инцидент.

Качество доставки вебхуков показывает доля уведомлений,
дошедших с~первой попытки; падение этой доли указывает либо на~проблемы
продавцов, либо на~перегрузку собственной очереди
доставки.

\section{REST API и~ISO~8583}\label{sec:gw-api-to-iso}

Пока разработчик смотрит на~транзакцию как на~JSON-объект
с~полями \texttt{amount}, \texttt{currency},
\path{payment_method} и~\texttt{description}, банк-эквайер
и~карточная сеть видят ту же операцию как сообщение
ISO~8583. В~публичных API современных PSP~--- международных
Stripe и~Adyen или местных ЮKassa, CloudPayments, T-Бизнес,
Сбер для~бизнеса~--- это буквально JSON поверх HTTPS с~чётко
описанным поведением запроса и~ответа~\cite{stripe-api-reference,yookassa-api,cloudpayments-api}.\footnote{Названия
эндпойнтов и~набор полей у~разных PSP свои. У~Stripe это
\texttt{POST~/v1/payment\_intents} и~\texttt{POST~/v1/charges},
у~ЮKassa~--- \texttt{POST~/v3/payments}, у~CloudPayments~---
\texttt{/payments/cards/charge} для~одностадийного платежа или
\texttt{/payments/cards/auth} для~двухстадийного, с~поддержкой
JSON и~form-encoded.
Иллюстрация ниже намеренно использует обобщённый
\texttt{POST~/charges} как нейтральный пример; дословным
эндпойнтом конкретного PSP она не~является.} Шлюз транслирует одно
представление в~другое; о~самом стандарте см.~главу~\ref{ch:iso8583}.

У~российских ТСП есть ещё один маршрут, не~имеющий аналога
в~карточных API международных PSP,~--- оплата через СБП
по~динамическому QR. На~уровне продукта это тот же вызов API вида
\texttt{POST~/payments} с~указанием способа оплаты СБП;
дальше PSP создаёт платёжную сессию у~банка-партнёра и~QR-код
в~формате СБП, а~шлюз возвращает продавцу ссылку или payload
для~отображения. Внутри маршрут идёт в~OpenAPI ОПКЦ СБП
(см.~раздел~\ref{sec:sbp-protocol}), минуя ISO~8583. Из-за этого
шлюз приходится проектировать как оркестратор методов оплаты
(компонент, координирующий несколько методов оплаты в~едином
сценарии), для~которого ISO~8583~--- лишь один из~исходящих
протоколов наряду с~OpenAPI ОПКЦ СБП и~собственными API
современных эквайеров.

Конвейер обработки одного \texttt{POST~/charges} внутри шлюза устроен так.

\begin{enumerate}
  \item \textbf{Приём и~валидация}~--- HTTP-слой проверяет
    ключ API, парсит JSON, валидирует обязательные поля
    (\texttt{amount}~>~0, \texttt{currency} по~ISO~4217,
    наличие \path{payment_method}). Если передан
    \texttt{idempotency-key}, шлюз ищет кэшированный ответ
    и~при совпадении возвращает его без повторной обработки.
  \item \textbf{Скоринг фрода}~--- шлюз вычисляет оценку мошенничества сам или
    запрашивает у~внешнего сервиса. Если оценка выше порога
    отказа, транзакция отклоняется до~контакта с~эквайером,
    экономя стоимость авторизационного запроса
    (см.~главу~\ref{ch:fraud}).
  \item \textbf{Маршрутизация}~--- модуль маршрутизации выбирает
    эквайера по~диапазону BIN, валюте, MCC и~текущей доле
    одобрений каждого канала. Если основной канал деградировал,
    каскадное правило переключает трафик на~резервный маршрут.
  \item \textbf{Извлечение PAN}~--- если \path{payment_method}
    ссылается на~сохранённый токен, шлюз обращается к~хранилищу
    vault за~FPAN (Funding PAN~--- реальный номер карты, в~отличие
    от~DPAN, токенизированного) или непосредственно за~DPAN. Хранилище
    возвращает карточные данные в~оперативную память процесса;
    на~диск они не~попадают
    (см.~раздел~\ref{sec:pci-scope-reduction}).
  \item \textbf{Сборка ISO~8583}~--- шлюз формирует сообщение
    с~индикатором типа MTI~0100 (Message Type Indicator),
    заполняет битовую карту присутствия полей (bitmap) и~поля
    данных (data elements, DE) согласно интерфейсной спецификации
    (IFS, Interface Specification) выбранного эквайера; вводный
    разбор формата дан в~главе~\ref{ch:iso8583}.
    Сопоставление полей описано ниже.
  \item \textbf{Отправка эквайеру}~--- сериализованное сообщение
    уходит по~соединению TCP/TLS из~пула. Шлюз запускает
    таймер ожидания ответа; верхний предел задают правила
    платёжной системы и~согласованное с~эквайером соглашение
    об~уровне обслуживания (SLA, Service-Level Agreement).
    Для~НСПК внутрироссийский авторизационный таймаут
    исчисляется десятками секунд; конкретный пункт сверяется
    по~действующей редакции Правил платёжной системы
    \enquote{Мир}~\cite{nspk-mir-rules}.
  \item \textbf{Парсинг ответа}~--- ответ MTI~0110 парсится,
    DE~39 (Response Code) транслируется в~статус JSON. Шлюз
    сохраняет результат и~обновляет запись транзакции.
  \item \textbf{Событие реестра}~--- шлюз публикует событие
    \texttt{authorization.created} во~внутреннюю шину. Модуль
    реестра обязательств фиксирует pending debit/credit
    (см.~раздел~\ref{sec:ledger-events}).
  \item \textbf{Вебхук}~--- событие ставится в~очередь доставки
    продавцу, подписывается ключом HMAC эндпоинта и~отправляется
    на~URL обратного вызова с~повтором при неуспехе.
  \item \textbf{Синхронный ответ}~--- HTTP-слой возвращает JSON
    продавцу с~полями \texttt{status}, \texttt{id},
    \path{decline_code} (при отказе) и~\texttt{timestamp}.
\end{enumerate}

Конкретные значения p50/p95 латентности конвейера зависят
от~множества факторов: расположения шлюза, эквайера и~эмитента,
протокола сети, нагрузки HSM, политики повторов. Публичных
агрегированных бенчмарков по~местному рынку нет, поэтому
числа в~книге не~приводятся; ориентир для~собственного
измерения~--- разделение полного времени ответа на~сегменты:
внутренняя обработка шлюза, ожидание эквайера, ожидание
сети и~эмитента.

Сборка ISO~8583 на~шаге~5 превращает короткий API-запрос в~сообщение
из~десятков обязательных полей. Продавец отправляет запрос с~суммой
\texttt{1500}, валютой \texttt{RUB}, номером карты и~сроком действия.
Шлюз транслирует это в~MTI~0100 (Authorization Request), заполняя:

\begin{itemize}
  \item \texttt{amount} 1500 $\to$ DE~4 Amount, Transaction~---
    12-символьное числовое поле, выровненное нулями,
    \texttt{000000150000} в~минорных единицах, копейках;
  \item \texttt{currency} \texttt{RUB} $\to$ DE~49 Currency Code,
    Transaction~--- \texttt{643}, числовой код рубля по~ISO~4217;
  \item \texttt{card.number} $\to$ DE~2 Primary Account Number~---
    PAN, до~19 цифр;
  \item \path{card.exp_month}, \path{card.exp_year} $\to$
    DE~14 Date, Expiration~--- \texttt{YYMM};
  \item внутренний ID терминала ТСП $\to$ DE~41 Card Acceptor
    Terminal Identification~--- 8-символьный код;
  \item MCC ТСП $\to$ DE~18 Merchant Type~--- 4-значный код.
\end{itemize}
Шлюз также генерирует поля, которых нет в~запросе API
продавца, но~требуемые протоколом~--- DE~11 (STAN), уникальный
в~коротком интервале, обычно сутки или смена, 6-значный
номер транзакции, DE~7 (Transmission Date
and Time), DE~12 (Time, Local Transaction), DE~22 (Point
of Service Entry Mode)~--- способ ввода данных карты
и~канал предъявления,~--- и~DE~25 (Point of Service Condition Code,
условия предъявления карты в~точке обслуживания); DE~22 кодирует,
\emph{как} были введены данные карты (магнитная полоса, чип,
бесконтакт, ручной ввод), а~DE~25~--- в~\emph{каких} условиях это
произошло (например, операция без присутствия держателя).

Если продавец использует сетевой токен вместо PAN, шлюз заполняет
DE~2 токенизированным PAN (DPAN), а~криптограмму от~TSP
передаёт в~EMV~data~--- данные транзакции по~стандарту EMV,
обычно переносимые в~DE~55,~--- или в~специфичном для~схемы поле,
зависящем от~конкретного протокола
маршрута~\cite{emvco-payment-tokenisation-guide}
(см.~также раздел~\ref{sec:token-anatomy}).

Ответ эмитента проходит обратный путь. DE~39 (Response Code),
двухсимвольный код, шлюз транслирует в~JSON-поле
\texttt{status} со~значениями \texttt{succeeded}, \texttt{declined},
\path{requires_action} и~дополняет человекочитаемым
описанием. Шлюз \emph{теряет} часть
информации: ISO~8583 различает десятки причин отказа
(\texttt{05}~Do Not Honor, \texttt{51}~Insufficient Funds,
\texttt{61}~Exceeds Withdrawal Limit), тогда как API шлюза группирует их
в~несколько категорий. Чтобы отличить
\enquote{недостаточно средств} от~\enquote{карта заблокирована},
разработчику нужно смотреть на~поле \path{decline_code}
или его аналог в~дополнение к~\texttt{status};
подробнее об~интерпретации кодов ответа см.~\cite{mc-tpr}
и~раздел~\ref{sec:decline-response-codes}.

\practicenote{%
  Многие шлюзы предоставляют подробный протокольный лог транзакции
  (raw response) со~значениями полей данных DE~39,
  DE~44 (Additional Response Data) и~DE~55 (EMV~data). Если шлюз
  этого не~показывает и~эквайер возвращает непонятный отказ,
  запросите у~эквайера trace log (низкоуровневую запись обмена
  ISO-сообщениями) по~STAN и~RRN~--- это единственный
  способ понять, что именно произошло на~уровне протокола.
}

\enlargethispage{4\baselineskip}
Описанная трансляция охватывает только авторизацию. Возврат, отмена
(void) и~подтверждение не~сводятся к~одному универсальному набору
сообщений для~всех шлюзов и~эквайеров: конкретное соответствие
зависит от~стадии операции и~протокола маршрута. На~уровне API
это несколько продуктовых действий, а~на~уровне карточного
обмена~--- последовательность авторизаций, отмен, финансовых
сообщений (financial messages, обычно MTI~0200/0220) и~клиринговых
обновлений (clearing updates), которые потом оказываются
во~внутреннем реестре обязательств и~в~сверке.

Авторизация укладывается в~одну пару MTI~0100/0110 (см.~раздел~\ref{sec:iso8583-mti}),
а~отмена, возврат и~подтверждение расходятся по~разным сообщениям
и~разным фазам жизненного цикла транзакции; различия reversal, void
и~refund разобраны в~разделе~\ref{sec:tx-cancellation}. Отмена авторизации
(reversal) уходит эквайеру до~клиринга~--- MTI~0400 в~интерфейсах
с~синхронным ответом эмитента, MTI~0420 в~интерфейсах с~односторонним
уведомлением (см.~раздел~\ref{sec:iso8583-mti}). Если эквайер не~отправил
reversal в~окне холда, заданном правилами сети (у~Mastercard~--- семь
суток для~final authorization и~до~тридцати для~pre-authorization),
эмитент сам снимает холд по~истечении срока, а~единственным способом
вернуть деньги после клиринга остаётся отдельный refund~\cite{mc-tpr}.
Возврат после клиринга в~single-message-сети идёт сообщением
MTI~0200 с~processing code~20 (refund), в~dual-message-сетях~--- отдельным
финансовым представлением через клиринг: IPM~1240 First Presentment
у~Mastercard, запись TC~33.A в~Capture File у~Visa~\cite{mc-gcms-rm-2021,visa-tc33a-capture-file}.
В~реестре обязательств refund появляется самостоятельной проводкой
со~ссылкой на~исходную авторизацию, не~как её зеркало
(см.~раздел~\ref{sec:tx-cancellation}). Подтверждение двухстадийной
авторизации (capture, финансовое представление) закрывает удержание
и~переводит сумму в~клиринг~--- MTI~0220 у~одних протоколов,
финансовое представление через клиринговый файл у~других, отдельный
REST-эндпойнт у~современных эквайеров~\cite{mc-gcms-rm-2021,visa-tc33a-capture-file}.
Шлюз сводит продуктовые \texttt{refund}, \texttt{void} и~\texttt{capture}
к~разным комбинациям сообщений на~каждом канале.
